\begin{statement}{E.33}{Proved}{The full rank-ten class-two refinement at three}
Retain the actual groups $H\triangleleft F\triangleleft B$,
the ten signed Schreier generators $\eta_1,\ldots,\eta_{10}$
of E.26, and the common Ferrari source $\mathscr D$ of E.24.
Put
\[
 P_3=B/\bigl(\gamma_3(H)H^3\bigr),\qquad
 \mathscr J_3=\operatorname{im}
       \bigl(B\longrightarrow\mathscr D\times P_3\bigr).
 \tag{E.33.1}
\]
The geometric kernel $H/(\gamma_3(H)H^3)$ has order $3^{55}$,
with ten abelian and forty-five commutator coordinates.
Every one of the $432$ multiplier-one points of E.24 has
a generating lift to the actual common source $\mathscr J_3$.
For each such point the complete set of correction images
coming from $\ker(\widehat B\to\mathscr D)\cap\widehat F'$
and satisfying both order relations has exactly
\[
                         3^{11}
 \tag{E.33.2}
\]
elements.

This refinement imposes a strict additional condition on
the abelian boundary fibre.  For a commutator word $f$,
write its $H$-abelianization modulo three as $v$.
The abelian order relations and the exact $F$-abelianization
condition give
\[
 v=v(a,b,c)=
 (c,b,a,b,a+2c,2a+c,2a,a+2c,c,2c)^t,
       \qquad a,b,c\in\mathbf F_3.
 \tag{E.33.3}
\]
The class-two order relations admit a central coordinate
over $v$ exactly when
\[
                       a-b-c=0,
 \qquad\text{equivalently}\qquad v_3-v_2-v_1=0.
 \tag{E.33.4}
\]
Each admissible $v$ has an affine ten-dimensional space of
central coordinates.  In each complete E.24 correction fibre,
exactly three of the nine abelian solutions of E.31 satisfy
\textup{(E.33.4)}.  Thus six actual abelian norm solutions in
each such fibre do not extend to \textup{(E.33.1)}, while all
$432$ source points retain generating extensions.
Every one of the $3^{11}$ finite solutions over a given
source point extends compatibly and simultaneously through
\[
 \mathscr J_{3,m}=
 \operatorname{im}\bigl(
 B\longrightarrow
 \mathscr D\times P_3\times B/[H,H]H^m\bigr),
                    \qquad m\ge1.
 \tag{E.33.17}
\]
One profinite commutator word realizes each such compatible
family, retaining all the Ferrari coordinates.
\end{statement}
\begin{proof}
The normality of the subgroup in \textup{(E.33.1)} follows
from characteristicness inside the actual normal subgroup
$H$.  For $q=(a,b)\in(\mathbb Z/3)^2$ use the section
$c(q)=x^ay^b$, with $0\leq a,b<3$.  The nontrivial signed
Schreier edges
$c(q)\ell c(q+\ell)^{-1}$, in the fixed order, are
\[
 \begin{gathered}
 ((0,1),x),\ ((0,2),x),\ ((0,2),y),\
 ((1,1),x),\ ((1,2),x),\ ((1,2),y),\\
 ((2,0),x),\ ((2,1),x),\ ((2,2),x),\ ((2,2),y).
 \end{gathered}
 \tag{E.33.5}
\]
This specifies every generator and every based comparison
without a choice of a further conjugating path.

Let $V=\mathbf F_3^{10}$, with this ordered basis, and
$Z=\Lambda^2V$.  The explicit isomorphism
\[
 H/\bigl(\gamma_3(H)H^3\bigr)\longrightarrow V\oplus Z,
 \qquad \eta_i\longmapsto(e_i,0),
 \tag{E.33.6}
\]
uses the multiplication
\[
 (v,z)(w,t)=\bigl(v+w,z+t+2v\wedge w\bigr).
 \tag{E.33.7}
\]
Hall collection gives the inverse map.  The ordinary
commutators are the independent coordinates of $Z$;
all cubes in \textup{(E.33.7)} are one.  Conversely the
class-two collection identities and
$\binom32=3$ show that the relations $\gamma_3(H)=H^3=1$
introduce no further relation among these coordinates.
This proves \textup{(E.33.6)}, rather than only a graded
dimension count.  Since $[B:H]=54$, the order of $P_3$
is $54\cdot3^{55}$.

Here are the exact action and norm calculations.  On the
free group $F$ the actual substitutions are
\[
 s:(x,y)\longmapsto(y,x),\qquad
 r:(x,y)\longmapsto\bigl(y,(xy)^{-1}\bigr).
 \tag{E.33.8}
\]
They are the negative-square convention verified in E.26.
Substitute each of the ten words \textup{(E.33.5)}, perform
the signed Schreier rewrite, and collect by
\textup{(E.33.7)}.  This gives matrices
\[
 \mathsf S=
 \begin{pmatrix}S_1&0\\J_s&\Lambda^2S_1\end{pmatrix},
 \qquad
 \mathsf R=
 \begin{pmatrix}R_1&0\\J_r&\Lambda^2R_1\end{pmatrix}.
 \tag{E.33.9}
\]
The computed matrices satisfy
$\mathsf S^2=\mathsf R^3=1$ and preserve the displayed
bracket.  Their Johnson blocks have ranks seven and nine.
The actual words $x=(r^{-1}s)^2$ and
$y=(sr^{-1})^2$ give precisely conjugation by $x$ and $y$
on all fifty-five coordinates.  Thus no diagonal action
has been substituted for the based one.

For completeness, the finite group operations used in the
certificate also have an explicit source comparison.  Write
an arbitrary $F$ word as $h\,c(q)$ by signed Schreier
rewriting.  Multiplication is
\[
 (h,q)(k,t)=
 \bigl(h\,{}^{c(q)}k\,
       c(q)c(t)c(q+t)^{-1},\,q+t\bigr).
 \tag{E.33.10}
\]
Both the conjugation and each of the eighty-one carries in
this formula are collected from their literal words.
This evaluates every relator used below in the actual
quotient, including all power terms.

For $g=(v,z)$ let multiplication in the following expression
mean \textup{(E.33.7)}, and let the displayed product be its
logarithmic coordinate:
\[
 \mathcal N(g)=
 \begin{pmatrix}
     (\mathsf Sg)\,g\\
     (-g)\,(\mathsf R(-g))\,(\mathsf R^2(-g))
 \end{pmatrix}.
 \tag{E.33.11}
\]
These are exactly $(sf)^2$ and $(f^{-1}r)^3$.  Their
abelian parts are
$(1+S_1)v$ and $-(1+R_1+R_1^2)v$.
The map $H^{\rm ab}\to F^{\rm ab}$ is the integral
map $3L$ of E.28.  A word in $F'$ therefore has
$L[f]=0$ before reduction, and hence $Lv=0$.
Solving these three literal linear conditions gives
\textup{(E.33.3)}.  The $Lv=0$ condition is retained;
the two norm conditions alone have a four-dimensional
kernel modulo three.

The central part of \textup{(E.33.11)} has the form
\[
                  Mz+e(v),\qquad
 M=\begin{pmatrix}
       1+\Lambda^2S_1\\
       -1-\Lambda^2R_1-(\Lambda^2R_1)^2
    \end{pmatrix}.
 \tag{E.33.12}
\]
Its rank is $35$, so its kernel has dimension ten.
There is a small certificate for the extra obstruction.
For the central coordinates $t_{ij}$ of the second
norm component put
\[
 \begin{split}
 \omega(t)={}&t_{13}+t_{14}-t_{15}+t_{16}+t_{17}-t_{18}\\
             &+t_{23}+t_{24}+t_{27}-t_{28}.
 \end{split}
 \tag{E.33.13}
\]
Literal matrix multiplication gives $\omega M=0$ and
\[
                 \omega e(v(a,b,c))=a-b-c .
 \tag{E.33.14}
\]
The certificate verifies this as a polynomial identity:
the three linear coefficients are $(1,-1,-1)$ and
every square and cross coefficient is zero.
It also solves \textup{(E.33.12)} for each of the nine
vectors $v(b+c,b,c)$.  These forty-five-coordinate
particular solutions, together with a basis of $\ker M$,
are saved in the exact certificate.  This proves both
directions of \textup{(E.33.4)}, and not only a necessary
linear test.

We next prove that the correction image is complete.
Put $R_3=\ker(B\to D_6)\cap F'$.
Use the complete $6,668$ normal relators and their actual
word metadata from E.31: the $4,656$ rank-$97$
commutators, the $95$ primitive integral-$F^{\rm ab}$
kernel power products, the eighteen basis cubes,
the $153$ basis commutators and the $1,746$
conjugation relations, each corrected by its literal
eighteen-coordinate representative.  Evaluate every
one by \textup{(E.33.10)}.  Their logarithms span a
$51$-dimensional subspace, already stable under
$\mathsf S,\mathsf R$ and every bracket.  Its abelian
projection is exactly the six-dimensional reduction
$\overline K$ of the integral lattice $K$ in E.31,
and its central intersection is all of $Z$.  Thus
\[
       \operatorname{im}(R_3\to V\oplus Z)
                         =\overline K\oplus Z.
 \tag{E.33.15}
\]
The logarithmic span equals the group image: in
\textup{(E.33.7)}, addition is obtained from group
multiplication by multiplying by the appropriate
half-power of the commutator, and scalar multiplication
is ordinary group powering.  Stability under both actual
$B$ generators and brackets therefore proves equality
with the full normal image.  Every chosen spanning
vector retains a relator or conjugation/commutator
preimage in the saved certificate.

The same image is obtained from
$\ker(B\to\mathscr D)\cap F'$.
Indeed the binary $K_6$ image has exponent two.
For each $u\in R_3$, the word $u^4$ kills that binary
image, kills the ternary image, and has the same
image as $u$ in the exponent-three group
\textup{(E.33.6)}.  This is an actual word map
retaining all induced Ferrari coordinates.
Density of the discrete commutator subgroup in
the full profinite commutator subgroup, intersected
with the specified open source kernel, gives the
same finite image for the profinite notation in
the statement.

For each retained literal word $f$, the allowed targets
therefore form exactly its coset by
$\overline K\oplus Z$.  On $\overline K$ the abelian
norm map has rank four and kernel dimension two,
giving the nine abelian options.  The literal
$432$-word reconstruction computes these affine
planes and verifies that each meets
\textup{(E.33.4)} in a line of three elements.
Each point of that line has $3^{10}$ central
solutions by \textup{(E.33.12)}.  Multiplication
by the fixed value of $f$ is a bijection on
correction images, proving the exact count
\textup{(E.33.2)}.
For example the actual E.24 point with indexed
CM row $2$, ternary parameters $(0,0)$ and binary
parameter $0$ has the abelian solution
\[
                 (0,0,1,0,1,2,2,1,0,0)^t.
 \tag{E.33.16}
\]
Its obstruction in \textup{(E.33.14)} is one.
This supplies a specified boundary solution that
does not lift, while the other three abelian
solutions for that source point do lift.

There is an exact kernel intersection that supplies the
simultaneous assertion.  Let $g_1,\ldots,g_8\in R_3$
be the actual words of E.31 whose integral
$H$-abelianizations are the basis columns of $K$.
For an original normal relator $u$ write, uniquely,
$[u]_H=Kc_u$ with $c_u\in\mathbb Z^8$, and form
\[
 u^\circ=u\left(\prod_{j=1}^8g_j^{(c_u)_j}\right)^{-1}.
 \tag{E.33.18}
\]
This is an actual element of $R_3\cap[H,H]$.
The program
\texttt{bc\_fermat\_class\_two\_integral\_kernel.py}
computes the exact integral $c_u$, using the same
eight literal words, and evaluates \textup{(E.33.18)}.
Forty-five of these adjusted relators have linearly
independent central images modulo three.  Their
indices and all eight integral exponents are saved.
Consequently
\[
 \operatorname{im}\bigl(R_3\cap[H,H]\longrightarrow Z\bigr)
                              =Z.
 \tag{E.33.19}
\]
This equality follows already from these forty-five
word preimages; it requires no assertion that a
finite kernel automatically lifts integrally.
Taking fourth powers again retains their central
images and puts them in
$\ker(B\to\mathscr D)\cap[H,H]$.

Now prescribe any one of the finite solutions counted
in \textup{(E.33.2)}.  E.31's proved finite-fibre
comparison lifts its abelian correction modulo three
through the entire three-adic norm fibre.  The two-adic
and prime-to-six supported corrections of that statement
give one word realizing the abelian tower at every
modulus with that specified reduction.
The central discrepancy from the prescribed $P_3$
solution belongs to $Z$.  Correct it by a product
of the actual fourth powers in
\textup{(E.33.18)--(E.33.19)}.
These words have zero integral $H$-abelianization,
so every abelian modulus and every Ferrari coordinate
is retained.  The resulting single word has the
prescribed $P_3$ image and both norm equations in
every factor of \textup{(E.33.17)}.

Finally all these finite norm solutions generate
the common source.  Let $W$ be the image of $N_B$
in $\mathscr J_3$.  Its image in $\mathscr D$
is finite nilpotent with prime divisors two and
three, and its image in $P_3$ is a three-group:
the latter maps to the cyclic-three image in
$F/H$ with a three-group kernel.
Consequently $W$ is finite nilpotent, is generated
by the $97$ specified generators, and maps onto
$(\mathbb Z/6)^{97}$ through the retained $G_6$.
Its two Sylow Frattini quotients therefore both
have dimension exactly $97$.
The corrected map is a homomorphism into the
actual common image, because both norm equations
hold in both factors.  It preserves $N_B$ through
its retained $Q$ map and induces the known
$G_6$ automorphism on these Frattini quotients.
The finite Frattini criterion gives surjectivity
on $W$, and then on $\mathscr J_3$.
This proves generation without imposing invariance
of the fixed Ferrari source kernel.

Generation in \textup{(E.33.17)} follows through the
same actual groups.  First take $m=2^a3^b$ with
$a,b\ge1$.  Choose a lower-central index larger
than the nilpotence classes of the images of $N_B$
in $\mathscr D$, in $P_3$ and in the three-power
abelian boundary quotient.  Its image in the
two-power boundary quotient is exactly
$\mathcal D/2^a\mathcal D$ by E.28.
The corrected image projects onto that
two-power boundary group, since the already
generating abelian action has not changed.
Thus this deep commutator subgroup of the corrected
joint image is exactly the pure two-primary factor
$\mathcal D/2^a\mathcal D$.
After quotienting it out, the joint $N_B$ image is
finite nilpotent with primes two and three and
has the same two rank-$97$ Frattini quotients
through $G_6$.  The preceding Frattini argument
therefore proves joint surjectivity.
For a further factor of $m$ prime to six, the
actual $H$ image in all existing factors is a
two-three group, whereas the additional abelian
$H$ image has order prime to six.  The exact
CRT map and coprime subgroup argument of E.31
give their full product over $B/H$.
The added $P_3$ factor changes neither coprimality
nor the prescribed identity action on this
prime-to-six boundary factor.  This proves
surjectivity for every $m$; the same constructed
word proves compatibility.

The programs
\texttt{bc\_fermat\_class\_two.py},
\texttt{bc\_fermat\_class\_two\_kernel.py},
\texttt{bc\_fermat\_class\_two\_fibres.py}, and
\texttt{bc\_fermat\_class\_two\_obstruction.py}
save the full matrices, literal preimage metadata,
all finite solutions and the polynomial obstruction
certificate.  Each executed output ends in
\texttt{PASS}.  The additional integral-kernel
certificate supplies the actual preimages needed
for the simultaneous all-moduli comparison.
\end{proof}
